% elementi raccolti dai primi capitoli delle linee guida e dal libro What is the text encoding initiative (Lou Burnard)
% what is the Text Encoding Initiative
% The TEI as an organization is committed to improving the accessibility and comprehensibility of the TEI as a technical system.

\documentclass{beamer}
    
%    \usepackage[english]{babel}
    %\usepackage[latin1]{inputenc}
    %\usepackage[T1]{fontenc}

\mode<presentation>{
  \setbeamertemplate{background canvas}[vertical shading]
  \usetheme{Berkeley}
  \useoutertheme{himinfolines}
}
  
\usepackage{ucs}
\usepackage[utf8]{inputenc}
\usepackage[english,polutonikogreek,italian,UKenglish,british]{babel}
\usepackage{graphicx}
\usepackage{colortbl}
\usepackage{multicol}
\usepackage{ulem}
\usepackage{verbatim}
\usepackage{alltt}
\usepackage{ccicons}
\usepackage{MnSymbol,wasysym}
\usepackage{tikzsymbols}
\usepackage{textcomp}
\usepackage{xmpincl}

\usepackage{parskip}
\setcounter{nframes}{100}
\setcounter{nframe}{1}
\setbeamercovered{dynamic}
\newenvironment{grcenv}{\begin{otherlanguage}{greek}}{\end{otherlanguage}}
\newcommand{\g}[1]{\textgreek{#1}}
\definecolor{darkgreen}{rgb}{0,0.5,0}
\definecolor{darkblue}{rgb}{0,0,0.5}
\definecolor{grey}{rgb}{0.5,0.5,0.5}
\setcounter{tocdepth}{5}

\makeatletter

\makeatother
%\includexmp{LicencesAndLicensing}

%frame00 metadata
    \title{Elementi Codifica TEI}
    \author[A.M. Del Grosso]{Angelo Mario Del Grosso}
    \institute{\texttt{angelo.delgrosso@ilc.cnr.it} \\\bigskip\textit{CNR-ILC-LicoLab} \\\bigskip\url{http://licolab.ilc.cnr.it/}}
    \date{Istituto di Linguistica Computazionale ``A. Zampolli'', \today}
    \AtBeginSection[]{
    \begin{frame}<beamer>
    \addtocounter{nframe}{1}
    \footnotesize
    \frametitle{Progress status}
    \tableofcontents[currentsection,hideothersubsections]
    \end{frame}
    }

\begin{document}

\begin{frame}
	\maketitle
\end{frame}

\begin{frame}
	\frametitle{Sommatio della Lezione}
	\tableofcontents
\end{frame}

\section{Intro}
\begin{frame}
	\frametitle{Introduzione}
	\addtocounter{nframe}{1}
    
    \begin{center}
	    \includegraphics[width=.2\textwidth]{../imgs/tei-r.pdf}
	\end{center}

    \begin{itemize}
        
        \item<1-> parleremo del sistema Modulare della TEI
            \begin{itemize}
                \item<1-> Moduli
                \item<1-> Classi
                \item<1-> Macro
                \item<1-> Datatype     
            \end{itemize} 
        \item<2-> parleremo degli elementi basilari
            \begin{itemize}
                \item<2-> Intestazione TEI (TEIHeader)
                \item<2-> Elementi e attributi globali
                \item<2-> Esempi di codifica     
            \end{itemize} 
    \end{itemize}
    
\end{frame}

\section{Sezione I}
%% sezione relativa alla infrestruttura TEI

% xx sezione 1 frame 01
\begin{frame}
    \frametitle{Attributi Globali}
    \framesubtitle{Elenco}
    \addtocounter{nframe}{1}


\textbf{\textrm{att.global} provides attributes common to all elements in the TEI encoding scheme.}

\begin{description}
    \item [@xml:id]     \textbf{identifier} provides a unique identifier for the element bearing the attribute.
    \item [@n]          \textbf{number} gives a number (or other label) for an element, which is not necessarily unique within the document.
    \item [@xml:lang]   \textbf{language} indicates the language of the element content using a ‘tag’ generated according to BCP 47\footnote{see \href{http://google.con}{http://google.com}}.
\end{description}

\end{frame}


% xx sezione 1 frame 02
\begin{frame}
    \frametitle{Attributi Globali}
    \framesubtitle{Elenco cont..}
    \addtocounter{nframe}{1}


\textbf{\textmd{att.global} provides attributes common to all elements in the TEI encoding scheme.}

\begin{description}
    \item rend [att.global.rendition]	(rendition) indicates how the element in question was rendered or presented in the source text.
    \item style [att.global.rendition]	contains an expression in some formal style definition language which defines the rendering or presentation used for this element in the source text
    \item rendition [att.global.rendition]	points to a description of the rendering or presentation used for this element in the source text.
\end{description}

\end{frame}

% xx sezione 1 frame 03
\begin{frame}
    \frametitle{Attributi Globali}
    \framesubtitle{Elenco cont...}
    \addtocounter{nframe}{1}


\textbf{\textmd{att.global} provides attributes common to all elements in the TEI encoding scheme.}

\begin{description}
    \item xml:base	provides a base URI reference with which applications can resolve relative URI references into absolute URI references.
    \item xml:space	signals an intention about how white space should be managed by applications.
    \item source [att.global.source]	specifies the source from which some aspect of this element is drawn.
\end{description}

\end{frame}

% xx sezione 1 frame 04
\begin{frame}
    \frametitle{Attributi Globali}
    \framesubtitle{Elenco cont....}
    \addtocounter{nframe}{1}


\textbf{\textmd{att.global} provides attributes common to all elements in the TEI encoding scheme.}

\begin{description}
    \item  cert [att.global.responsibility]	(certainty) signifies the degree of certainty associated with the intervention or interpretation.
    \item resp [att.global.responsibility]	(responsible party) indicates the agency responsible for the intervention or interpretation, for example an editor or transcriber.
\end{description}

\end{frame}


% xx sezione 1 frame 05

\begin{frame} [fragile]
    \frametitle{Attributi Globali}
    \framesubtitle{Esempio \textrm{@xml:lang}}
    \addtocounter{nframe}{1}

    \textbf{\textrm{xml:lang} indica la lingua e il sistema di scrittura usato}
    \defverbatim{\langatt}{%
        \begin{tiny}
        \begin{verbatim}
            <TEI xmlns="http://www.tei-c.org/ns/1.0">
                <teiHeader xml:lang="en">
                    <!-- ... -->
                </teiHeader>
                <text xml:lang="fr">
                    <body>
                        <div>
                            <!-- chapter one is in French -->
                        </div>
                        <div xml:lang="de">
                            <!-- chapter two is in German -->
                        </div>
                        <div>
                            <!-- chapter three is French -->
                        </div>
                        <!-- ... -->
                    </body>
                </text>
            </TEI>
        \end{verbatim}
        \end{tiny}
        }
        \begin{block}{XML TEI con esempio di uso \textrm{xml:lang} attribute}
            {\langatt}
        \end{block}
\end{frame}


% xx sezione 1 frame 06
\begin{frame}
    \frametitle{Attributi Globali}
    \framesubtitle{Stile e Aspetto}
    \addtocounter{nframe}{1}

    
    \textbf{In the TEI scheme, it is possible to supply information about the appearance of elements within a source document in the following distinct ways:}

    \begin{itemize}
        \item One or more properties may be specified as the default for a set of elements (based on an external scheme, by default CSS), using rendition elements and their selector attributes;
        \item One or more properties may be specified for individual element occurrences, using the rend attribute with any convenient set of one or more sequence-indeterminate tokens;
    \end{itemize}
\end{frame}

% xx sezione 1 frame 07
\begin{frame}
    \frametitle{Attributi Globali}
    \framesubtitle{Stile e Aspetto cont..}
    \addtocounter{nframe}{1}
    
    \textbf{Note that these TEI attributes always describe the rendition or appearance of the source document, not intended output renditions, although often the two may be closely related.}

    \begin{itemize}
        \item One or more properties may be specified for individual element occurrences, using the rendition attribute to point to rendition elements;
        \item One or more properties may be supplied explicitly for individual element occurrences, using the style attribute.
    \end{itemize}

\end{frame}


% xx sezione 1 frame 08

\begin{frame}
    \frametitle{Attributi Globali}
    \framesubtitle{da vari altri Moduli Tabella}
    \addtocounter{nframe}{1}
    %fare una tabella 
class name	module name	see further
att.global.linking	linking	16 Linking, Segmentation, and Alignment
att.global.analytic	analysis	17 Simple Analytic Mechanisms
att.global.facs	transcr	11.1 Digital Facsimiles
att.global.change	transcr	11.7 Identifying Changes and Revisions

\end{frame}


\section{Sezione II}
%% frame 01
% Every TEI document must have a TEI Header, represented by a <teiHeader> element.
% The TEI header has four main components,
% file description, encoding description, profile description, revision description

% esempio TEI header Minimale
% <teiHeader>
%    <fileDesc>
%        <titleStmt>
%            <title>Title of the work</title>
%        </titleStmt>
%        <publicationStmt>
%            <p>Information about the publication of the work</p>
%        </publicationStmt>
%        <sourceDesc>
%            <p>Information about the source from which the work was derived</p>
%        </sourceDesc>
%    </fileDesc>
% </teiHeader>

% Rifarsi al capitolo due delle linee guida per una descrizione esaustiva e sistematica

% esempio TEI descrizione file
% tre componenti obbligatorie: title statement, publication statement, descrizione della fonte.

% <titleStmt>
%    <title xml:lang="sk">Yogadarśanam (arthāt yogasūtrapūphah).</title>
%    <title>The Yoga sūtras of Patañjali: a digital edition.</title>
%    <author>Patañjali</author>
%    <funder>Wellcome Institute for the History of Medicine</funder>
%    <principal>Dominik Wujastyk</principal>
%    <respStmt>
%        <name>Wieslaw Mical</name>
%        <resp>data entry and proof</resp>
%    </respStmt>
%    <respStmt> 
%        <name>Jan Hajic</name> 
%        <resp>conversion to TEI-conformant markup</resp> 
%    </respStmt>
% </titleStmt>

% esempio publication statement

%% Descrizione della fonte
% document formally the object or objects from which the TEI document has been derived, using traditional bibliographic
% Born Digital
% Printed Source
% sbobbinatura
% manoscritto (msDesc)

%% Esempio di sourceDesc
% <sourceDesc>
%    <bibl xml:id="Sue1846">
%        <author>
%            <surname>Sue</surname>,
%            <forename>Eugène</forename>
%        </author>
%        <title level="m">Martin, l’enfant trouvé : Mémoires d’un valet de chambre</title>
%        <imprint>
%            <publisher>C. Muquardt</publisher>
%            <pubPlace>Bruxelles</pubPlace>
%            <pubPlace>Leipzig</pubPlace>
%            <date when="1846">MDCCCXLVI</date>
%        </imprint>
%    </bibl>
% </sourceDesc>


% file description 
\subsection{File Description}

 esempio TEI descrizione file
% tre componenti obbligatorie: title statement, publication statement, descrizione della fonte.

% <titleStmt>
%    <title xml:lang="sk">Yogadarśanam (arthāt yogasūtrapūphah).</title>
%    <title>The Yoga sūtras of Patañjali: a digital edition.</title>
%    <author>Patañjali</author>
%    <funder>Wellcome Institute for the History of Medicine</funder>
%    <principal>Dominik Wujastyk</principal>
%    <respStmt>
%        <name>Wieslaw Mical</name>
%        <resp>data entry and proof</resp>
%    </respStmt>
%    <respStmt> 
%        <name>Jan Hajic</name> 
%        <resp>conversion to TEI-conformant markup</resp> 
%    </respStmt>
% </titleStmt>

% esempio publication statement

%% Descrizione della fonte
% document formally the object or objects from which the TEI document has been derived, using traditional bibliographic
% Born Digital
% Printed Source
% sbobbinatura
% manoscritto (msDesc)

%% Esempio di sourceDesc
% <sourceDesc>
%    <bibl xml:id="Sue1846">
%        <author>
%            <surname>Sue</surname>,
%            <forename>Eugène</forename>
%        </author>
%        <title level="m">Martin, l’enfant trouvé : Mémoires d’un valet de chambre</title>
%        <imprint>
%            <publisher>C. Muquardt</publisher>
%            <pubPlace>Bruxelles</pubPlace>
%            <pubPlace>Leipzig</pubPlace>
%            <date when="1846">MDCCCXLVI</date>
%        </imprint>
%    </bibl>
% </sourceDesc>

% encoding description
\subsection{Encoding Description}
% Encoding descsription (<encodingDesc />) is used to supply information about almost any aspect of the encoding process itself

%% esempio descrizione della codifica
% <encodingDesc>
%    <projectDesc>
%        <p>Texts collected for use in the Claremont Shakespeare Clinic, June 1990.</p>
%    </projectDesc>
%    <samplingDecl>
%        <p>Each text contains a sample of up to 2000 words, running from the start of the document to the end of the sentence
%            after the 2000 word mark. For the purposes of word counting, hyphens and apostrophes were treated as spaces.
%            </p>
%    </samplingDecl>
%    <editorialDecl>
%        <normalization>
%            <p>Word forms broken by end of line hyphenation have been reconstructed without comment. The hyphen has been removed
%                except for hyphenated forms attested elsewhere in the text. </p>
%        </normalization>
%        <quotation marks="all" form="std">
%            <p>All quotation marks have been removed. Direct speech is represented by the use of the
%                <gi>said</gi> tag; other quoted material is represented by means of the
%                <gi>q</gi> tag. </p>
%        </quotation>
%    </editorialDecl>
% </encodingDesc>

%% for human reader and for automated process

% esempio charDecl
% <charDecl>
%    <glyph xml:id="z103">
%        <glyphName>LATIN LETTER Z WITH TWO STROKES</glyphName>
%        <mapping type="standardized">z</mapping>
%        <mapping type="PUA">U+E304</mapping>
%    </glyph>
% </charDecl>

% segue nel testo 
% <p> ... mulct<g ref="#z103"/> ... </p>

% tassonomia

% <classDecl>
%    <taxonomy xml:id="size">
%        <category xml:id="large">
%            <catDesc>story occupies more than half a page</catDesc>
%        </category>
%        <category xml:id="medium">
%            <catDesc>story occupies between quarter and a half page</catDesc>
%        </category>
%        <category xml:id="small">
%            <catDesc>story occupies less than a quarter page</catDesc>
%        </category>
%        <!-- etc -->
%    </taxonomy>
%    <taxonomy xml:id="topic">
%        <category xml:id="politics-domestic">
%            <catDesc>Refers to domestic political events</catDesc>
%        </category>
%        <category xml:id="politics-foreign">
%            <catDesc>Refers to foreign political events</catDesc>
%        </category>
%        <category xml:id="social-women">
%            <catDesc>refers to role of women in society</catDesc>
%        </category>
%        <category xml:id="social-servants">
%            <catDesc>refers to role of servants in society</catDesc>
%        </category>
%        <!-- etc -->
%    </taxonomy>
% </classDecl>

% uso nel testo 
% <catRef target="#small #social-women"/>

%% esempio dichiarazione dei tag
% <tagsDecl>
%    <rendition xml:id="IT" scheme="css">font-style:italic</rendition>
%    <rendition xml:id="FontRoman" scheme="css">font-family: serif</rendition>
%    <namespace name="http://www.tei-c.org/ns/1.0">
%        <tagUsage gi="emph" render="#IT" />
%        <tagUsage gi="hi" render="#IT" />
%        <tagUsage gi="text" render="#FontRoman" />
%    </namespace>
% </tagsDecl>

% uso nel testo
% @rendition attribute

% profile description
\subsection{Profile Description}
% elemento che descrive il profilo del documento TEI
% elementi di metadatazione non bibliografici
% elementi per edentificare i vari passaggi genetici elemento <listChange />
 %% esempio proust di Pierazzo
 %% esempio carteggio

 %% esempio profile
% <profileDesc>
%    <creation>
%        <date when="1962" /> </creation>
%    <textClass>
%        <catRef target="#WRI #ALLTIM1 #ALLAVA2 #ALLTYP3 #WRIDOM5 #WRILEV2 #WRIMED1 #WRIPP5 #WRISAM3 #WRISTA2 #WRITAS0" />
%        <classCode scheme="DLEE">W nonAc: humanities arts</classCode>
%        <keywords scheme="COPAC">
%            <term>History, Modern - 19th century</term>
%            <term>Capitalism - History - 19th century</term>
%            <term>World, 1848-1875</term>
%        </keywords>
%    </textClass>
% </profileDesc>

% revision description
\subsection{Revision Description}
% represented by a <revisionDesc> element which contains a list of <change>
% given first. The <listChange> element mentioned above may also be used here to refer to identified stages in the evolution of the electronic file

% esempio
% <revisionDesc>
%    <listChange>
%        <change when="2013-05-11">First complete draft</change>
%        <change when="2013-04-07">Created header and document structure</change>
%    </listChange>
% </revisionDesc>

\section{Sezione III}
% capitoli 3 e 4 delle linee guida e estratti dal libro what is TEI (The structural organization, Varieties of textual structure, parte della TEI cornucopia part 1)
% fare esempio facsimile e codifica dei fenomeni (vedere anche slide del turco e fiormonte-ciotti)


\section{Sezione IV}
% ODD document, Selezione dei Moduli per lo schema, nuovi elementi, profilo personalizzato TEI, capitolo 22 delle linee guida (Documentation Elements), capitolo 23 delle linee guida (Using the TEI)


\section{Conclusioni}
%conclusioni f01
\begin{frame}
    \frametitle{References}
    \addtocounter{nframe}{1}
    \begin{thebibliography}{10}

        \setbeamertemplate{bibliography item}[paper]
        \tiny\bibitem{ciotti2011} Ciotti 2011

        \setbeamertemplate{bibliography item}[online]
        \tiny\bibitem{CCwikiBY2014} CC Wiki \textit{Best practices for attribution}, CC Wiki 2014, \url{https://wiki.creativecommons.org/wiki/Best\_practices\_for\_attribution}

    \end{thebibliography}

\end{frame}


\end{document}